% Pending independent integration audit. Draft completed 2026-11-19.
\begin{saresult}{A.134}{Proved under stated hypotheses}{Real PL--TOP comparison for all closed bases of dimension at most two}
Let $B$ be a compact real PL manifold without boundary, all of whose
components have dimension at most two. The natural map from the simplicial
group of jointly PL concordances to the singular complex of the compact-open
topological concordance group is a weak equivalence:
\[
 P_{\mathrm{PL}}(B)\longrightarrow
 \operatorname{Sing}P_{\mathrm{TOP}}(B).
                                                        \tag{A.134.1}
\]
Both groups fix $B\times\{0\}$ pointwise and leave the top free.
For components of dimension zero or one these spaces are weakly
contractible. For every closed PL surface $F$, the natural comparison
$\operatorname{Aut}_{\mathrm{PL}}(F)\to
\operatorname{Sing}\operatorname{Homeo}(F)$ is a weak equivalence as well.
Together with A.131, this gives the explicit recursive whole-space TOP
model and recursive homotopy-group presentations in these base dimensions.
The empty base has the one-element simplicial group. No statement about
base dimension four, or about higher PL--semialgebraic comparison, is
included here.
\end{saresult}
\begin{proof}
\textbf{The exact joint-family input and its component convention.}
For a compact PL manifold $N$, write $\mathcal H_{\mathrm{PL}}(N)$ for
all projection-preserving PL homeomorphisms of $\Delta^q\times N$ in
simplicial degree $q$, and write $\mathcal H_{\mathrm{TOP}}(N)$ for the
singular complex of its entire compact-open homeomorphism group.
Hamstrom's boundary convention includes compact manifolds with boundary
\cite[\S0.2, p.~208]{Hamstrom1974}. Her \S2.1, p.~222, explicitly
separates the topology on the set of PL homeomorphisms from the
semisimplicial complex whose total maps $\Delta^q\times N$ are PL and
preserve the projection. Her homotopies preserve both simplex and homotopy
parameters. It is the latter, joint-family convention that her
Theorem~3.1.2, pp.~225--226, compares to TOP. The natural inclusion induces
isomorphisms on all positive homotopy groups when $\dim N=3$.
Adding the usual degeneracy operators does not change these homotopy
groups: the representative maps and relative homotopies just described
are precisely those of the stipulated simplicial group.

There is a convention to address before using this on all components.
Hamstrom's $\operatorname{Top}(N)$ means the TOP identity component, and
her PL vertices are the PL members of that component. We give the
component argument explicitly. The Moise--Bing approximation theorem
approximates every homeomorphism of a compact PL three-manifold by PL
homeomorphisms; see also its all-homeomorphism formulation in
\cite[Chapter~I, \S1, Theorem~1.1, p.~7]{Laudenbach1974}.
TOP local contractibility \cite{EdwardsKirby1971} implies that sufficiently
close homeomorphisms are TOP-isotopic. Consequently every TOP component
contains a PL vertex.

For injectivity on components, take a TOP isotopy $h_t$ between two PL
vertices. First right-multiply the whole isotopy by $h_0^{-1}$, so that
it starts at the identity and every member lies in the TOP identity
component. At the end right-multiply the constructed PL isotopy by $h_0$.
The joint PL local $0$-connectedness in
\cite[Theorem~2.3.1, p.~224, with the definition on p.~222]{Hamstrom1974}
supplies a $\delta>0$ such that every PL map $v$ with
$d(v,1)<\delta$ has a jointly PL isotopy to the identity. Subdivide the
compact parameter interval finely, and choose PL approximations $g_i$ to
$h_{t_i}$, retaining the prescribed PL endpoints, so that
$d(g_{i+1},g_i)<\delta$. For the uniform metric,
\[
 d(g_{i+1}g_i^{-1},1)=d(g_{i+1},g_i).
\]
The local PL isotopies, followed by right multiplication by $g_i$, join
successive $g_i$. Their finite concatenation is jointly PL. This proves
injectivity on $\pi_0$, without assigning a block meaning to an isotopy.
The approximation vertices lie in the TOP identity component whenever
the original path does; this also verifies the source convention used
in the local theorem. Translating by a PL representative now extends
Hamstrom's positive homotopy-group comparison to every component.
We have therefore proved the full, natural weak equivalence
\[
 \mathcal H_{\mathrm{PL}}(N)\longrightarrow
 \mathcal H_{\mathrm{TOP}}(N),\qquad \dim N=3.
                                                        \tag{A.134.2}
\]
Laudenbach's explicit ordinary-isotopy statement on p.~8 is consistent
with this component argument. His earlier footnote about a
quasi-simplicial complex is not used for the higher homotopy groups.

\textbf{A product decomposition for closed surfaces.}
Let $F$ be connected and closed, with $\dim F=2$, and let $N=F\times I$.
For either category let $\mathcal A(F)$ be the subgroup of
$\mathcal H(N)$ preserving $F\times\{0\}$ setwise.
The permutation of the two boundary components is constant on every
simplex and every path. Hence $\mathcal A(F)$ is a union of whole
components. The bijection on components and the isomorphisms at every
basepoint in (A.134.2) restrict to a weak equivalence
\[
 \mathcal A_{\mathrm{PL}}(F)\longrightarrow
 \mathcal A_{\mathrm{TOP}}(F).
                                                        \tag{A.134.3}
\]
Bottom restriction $r$ has the product section
$s(a)(x,u)=(a(x),u)$. There is an actual isomorphism of underlying
simplicial sets (and, in TOP, of the corresponding topological spaces)
\[
 \begin{split}
 \mathcal A(F)&\longrightarrow P(F)\times\mathcal H(F),\\
 h&\longmapsto\bigl(h\,s(rh)^{-1},rh\bigr),\\
 (p,a)&\longmapsto p\,s(a).
 \end{split}                                            \tag{A.134.4}
\]
Every formula preserves all parameters. Total PL maps have total PL
inverses, so the formula is jointly PL; in TOP composition and inversion
are continuous for the compact domain. The group structure is a
semidirect product, but only the literal product of underlying spaces
is needed.

Under (A.134.4), the natural comparison (A.134.3) is the product of the
concordance comparison and the homeomorphism comparison for $F$.
A product of maps between nonempty spaces can be a weak equivalence
only if each factor is: the assertion follows for components from the
bijection on product sets, and for every higher homotopy group from
the product decomposition at each basepoint. This proves (A.134.1)
for every closed connected surface, orientable or not, and proves the
assertion about $\mathcal H(F)$. No boundary-restriction fibration or
prescribed smoothing of $F$ is being assumed.

\textbf{The circle.}
We give a separate low-dimensional argument, rather than extrapolating
(A.134.2) to a manifold of the wrong dimension. Scott's original thesis
uses exactly the joint convention: Chapter~III, Definition~0.1, p.~19,
defines simplices to be PL homeomorphisms of $\Delta^q\times X$
commuting with projection \cite{Scott1968}. In Chapter~IV,
Theorem~3.1(b)(iii), p.~51, with the annulus argument on p.~53, he proves
that evaluation at an interior point induces isomorphisms
\[
 \pi_j\bigl(\mathcal H_{\mathrm{PL}}(S^1\times I),1\bigr)
 \longrightarrow\pi_j(S^1\times I),\qquad j\geq1.
                                                        \tag{A.134.5}
\]
Move the evaluation point to the bottom circle along a fixed path in
the annulus. Evaluation along that path gives a homotopy, so, with the usual basepoint transport, the same
assertion holds for the map given by bottom evaluation and inclusion
$S^1\times\{0\}\subset S^1\times I$.
Hamstrom's Theorem~2.2.4 and its proof, pp.~223--224, identify the
natural PL rotation circle with $\mathcal H_{\mathrm{PL}}(S^1)_0$;
evaluation on this circle has degree one. Thus evaluation at a point
induces the usual isomorphisms
$\pi_j(\mathcal H_{\mathrm{PL}}(S^1),1)\to\pi_j(S^1)$.
It follows from (A.134.5) that bottom restriction induces isomorphisms
on every positive homotopy group. Apply the same product formula
(A.134.4), now for the annulus. Its first factor has zero positive
homotopy groups. Its component group is zero by the circle argument
of A.123. Therefore $P_{\mathrm{PL}}(S^1)$ is weakly contractible.

For TOP, Yagasaki's compact-surface lemma states that
$\operatorname{Homeo}(S^1\times I\text{ rel }S^1\times\{0\})_0$
is contractible
\cite[Lemma~2.1(ii), p.~2, proof p.~3]{Yagasaki2000}.
We must still check that this is the whole group. If $h$ is a TOP
annulus homeomorphism fixed on the bottom, Epstein's Theorem~A4,
p.~107, isotopes $h$ to a PL annulus homeomorphism $g$
\cite{Epstein1966}. Throughout the isotopy $h_t$, the bottom boundary
component stays the bottom component, since a continuous family cannot
change its boundary permutation. Put $a_t=h_t|_{S^1\times\{0\}}$.
The corrected isotopy
\[
 k_t=(a_t^{-1}\times\operatorname{id}_I)\,h_t
                                                        \tag{A.134.6}
\]
is fixed on the bottom, starts at $h$, and ends at the PL concordance
$(a_1^{-1}\times\operatorname{id}_I)g$. A.123 joins this last map to
the identity through jointly PL concordances, hence also through TOP
concordances. Thus $P_{\mathrm{TOP}}(S^1)$ is connected and is its
contractible identity component. The map between these two weakly
contractible spaces is a weak equivalence.

\textbf{Points, disconnected bases, and the empty base.}
For a point, a concordance is an interval homeomorphism fixed at both
endpoints (fixing the bottom endpoint forces the other endpoint to be
fixed). The TOP group contracts by linear interpolation with the
identity. The joint PL simplicial group is weakly contractible by the
parameter-preserving Alexander construction
\cite[Theorem~2.2.1, p.~223]{Hamstrom1974}.
Equivalently, for a based finite simplex family, cone its total graph
from the barycenter of the parameter simplex and the center of the
interval, as in that theorem; this preserves the parameter projection
and is jointly PL. This formulation avoids falsely using the bilinear
interpolation formula as a jointly PL homotopy.

A compact manifold has finitely many connected components. A
concordance fixed on its entire bottom preserves each corresponding
cylinder component, so its simplicial group is the finite product of
the component groups, even if their dimensions differ. Finite products
preserve the proved comparisons. If there are no components, both
groups consist of the identity alone. Finally A.131 supplies the
explicit whole PL model and its recursive presentations, so their
transfer to TOP in the stated dimensions follows from (A.134.1).
\end{proof}
